\SGASubsection{C.2}{Complement to \SGARefLink{sga3:xv:remark:4:8}{Exposé XV, 4.8}}
\SGARefTarget{sga3:xvii:appendix:C:subsection:2}

The following proposition, as well as \SGARefLink{sga3:xvii:theorem:c-3-1}{Theorems C.3.1} and \SGARefLink{sga3:xvii:theorem:c-4-1}{C.4.1}
below, will appear in a forthcoming article by M. Raynaud on
group schemes over a discrete valuation ring.\footnote{Editor's
note: comments are to be added here, including references to
Exposé VI\textsubscript{B}, \ldots}

\SGAStatement{\SGAInlineRefTarget{sga3:xvii:proposition:c-2-1}Proposition C.2.1}{Let \(S\) be the spectrum of a
discrete valuation ring, let \(t\) be its generic point, let \(G\)
be a flat \(S\)-group prescheme of finite type, and put
\[
  \widetilde G=\operatorname{Spec}\Gamma(G,\mathcal O_G).
\]
Let \(u:G\to\widetilde G\) be the canonical morphism.  Then:
\begin{enumerate}
\item[(1)] \SGARefTarget{sga3:xvii:proposition:C:2:1:item:1}
  \(\widetilde G\) is naturally an \(S\)-group scheme, and \(u\)
  is a homomorphism.
\item[(2)] \SGARefTarget{sga3:xvii:proposition:C:2:1:item:2}
  \(\ker(u)\) is flat over \(S\), and
  \((\widetilde G,u)\) is the fpqc quotient of \(G\) by
  \(\ker(u)\).  Thus \(\widetilde G\) is the largest affine
  quotient of \(G\).
\item[(3)] \SGARefTarget{sga3:xvii:proposition:C:2:1:item:3}
  If \(G_t\) is affine, then \(\ker(u)\) is an étale group over
  \(S\), and it is the trivial group if and only if \(G\) is
  separated over \(S\).  In particular, a flat, separated
  \(S\)-group scheme of finite type whose generic fiber is affine
  is affine.
\end{enumerate}}
\SGARefTarget{sga3:xvii:proposition:C:2:1}

\SGANextSectionTarget{sga3:xvii:subsection:c-3}\SGASubsection{C.3}{}
\SGARefTarget{sga3:xvii:appendix:C:subsection:3}

In the statement of \SGARefLink{sga3:xv:lemma:6:6}{Exposé XV, 6.6}, the hypothesis that \(H\)
be equal to its connected normalizer is unnecessary in order that the
set of points \(s\) of \(S\) for which \(H_s\) is a parabolic
subgroup of \(G_s\) be open.  Indeed, repeat the proof given in
paragraph~c) preceding \SGARefLink{sga3:xv:lemma:6:9}{Lemma 6.9}, and let \(N\) denote the
scheme-theoretic closure of \(N_t\) in \(G\).  Then \(N\) is a
closed subgroup scheme of \(G\), flat over \(S\), which contains
\(H\) and is contained in \(\mathcal N\).  The theorem below
shows that \(G/N\) is representable.  Since \(G_s/N_s\) is proper
and connected, the proof concludes as in \emph{loc.\ cit.}

\SGAStatement{\SGAInlineRefTarget{sga3:xvii:theorem:c-3-1}Theorem C.3.1\footnote{Editor's note: see \SGARefLink{sga3:xvii:appendix:C:subsection:4}{Theorem
4}C in S. Anantharaman, ``Schémas en groupes, espaces homogènes et
espaces algébriques sur une base de dimension 1,'' \emph{Bull.
Soc. Math. France, Mém.} 33 (1976), 5--79.}}{Let \(S\) be the
spectrum of a henselian discrete valuation ring, let \(G\) be an
\(S\)-group prescheme locally of finite type, and let \(H\) be a
closed subgroup prescheme of \(G\), flat over \(S\).  Then \(G/H\)
is representable.}
\SGARefTarget{sga3:xvii:theorem:C:3:1}

\SGANextSectionTarget{sga3:xvii:subsection:c-4}\SGASubsection{C.4}{Complement to \SGARefLink{sga3:xvi:subsection:1:1}{Exposé XVI, 1.1}}
\SGARefTarget{sga3:xvii:appendix:C:subsection:4}

The following theorem refines \SGARefLink{sga3:viii:proposition:7:9}{Exposé VIII, 7.9}.

\SGAStatement{\SGAInlineRefTarget{sga3:xvii:theorem:c-4-1}Theorem C.4.1}{Let \(S\) be the spectrum of a
discrete valuation ring with residue characteristic \(p>0\), and
let \(G\) be a commutative, smooth, separated \(S\)-group scheme
of finite type.  Then the following conditions are equivalent:
\begin{enumerate}
\item[(i)] \SGARefTarget{sga3:xvii:theorem:C:4:1:item:i}
  \({}_pG\) is finite over \(S\).
\item[(ii)] \SGARefTarget{sga3:xvii:theorem:C:4:1:item:ii}
  For every \(S\)-prescheme \(S'\), every separated
  \(S'\)-group prescheme \(H'\) of finite presentation over
  \(S'\), and every \(S'\)-monomorphism
  \[
    u:G_{S'}\longrightarrow H',
  \]
  the morphism \(u\) is an immersion.
\end{enumerate}}
\SGARefTarget{sga3:xvii:theorem:C:4:1}

\SGANextSectionTarget{sga3:xvii:subsection:c-5}\SGASubsection{C.5}{}
\SGARefTarget{sga3:xvii:appendix:C:subsection:5}

Example XVII, \SGARefLink{sga3:xvii:subsection:6:4:item:a}{6.4(a)}, gives a smooth group over a field \(k\)
whose unipotent radical is not defined over \(k\).  The following
proposition gives a general method for obtaining such groups.

\SGAStatement{\SGAInlineRefTarget{sga3:xvii:proposition:c-5-1}Proposition C.5.1}{Let \(k\) be a field, let \(K\)
be a finite radicial extension of \(k\) of degree greater than
\(1\), and let \(H\) be a smooth, connected \(K\)-algebraic group
of dimension \(r\).  Put
\[
  G=\prod_{K/k}H/K,
\]
which is a smooth, connected \(k\)-algebraic group of dimension
\(r[K:k]\).  Let
\[
  u:G_K\longrightarrow H
\]
be the canonical homomorphism, and put \(R=\ker(u)\).  Then:
\begin{enumerate}
\item[(i)] \SGARefTarget{sga3:xvii:proposition:C:5:1:item:i}
  \(u\) is an epimorphism, and \(R\) is a smooth, connected,
  unipotent algebraic group.
\item[(ii)] \SGARefTarget{sga3:xvii:proposition:C:5:1:item:ii}
  If \(U\) is a smooth algebraic subgroup of \(G\) such that
  \(U_K\) contains \(R\), then \(U=G\).
\end{enumerate}}
\SGARefTarget{sga3:xvii:proposition:C:5:1}

\SGAStatement{Corollary}{Retain the notation of the preceding
proposition.
\begin{enumerate}
\item[(a)] \SGARefTarget{sga3:xvii:corollary:C:5:1:item:a}
  If \(H\) is not unipotent, then the unipotent radical of
  \(G_{\bar k}\) is not defined over \(k\).
\item[(b)] \SGARefTarget{sga3:xvii:corollary:C:5:1:item:b}
  If \(H\) is a nonzero abelian variety, then \(G\) is not an
  extension of an abelian variety by a smooth linear group.
\item[(c)] \SGARefTarget{sga3:xvii:corollary:C:5:1:item:c}
  If \(H\) is not solvable, then the solvable radical of
  \(G_{\bar k}\) is not defined over \(k\), and \(G\) has no
  Borel subgroup defined over \(k\).
\end{enumerate}}
\SGARefTarget{sga3:xvii:corollary:C:5:1}

Let us first show how the corollary follows from the proposition.

\smallskip
\noindent\textit{\SGARefLink{sga3:xvii:appendix:A}{(a)}} Let \(U\) be a unipotent radical of \(G\).
Then \(U_K\) is the unipotent radical of \(G_K\), and hence
contains \(R\), since \(R\) is smooth, connected, and unipotent
by \SGARefLink{sga3:xvii:corollary:C:5:1:item:i}{(i)}.  Thus \(U=G\) by \SGARefLink{sga3:xvii:corollary:C:5:1:item:ii}{(ii)}.  But \(G_K\) has \(H\) as a
quotient, and \(H\) is not unipotent by hypothesis.  Therefore
\(G\) is not unipotent, a contradiction.

\smallskip
\noindent\textit{\SGARefLink{sga3:xvii:section:b}{(b)}} Suppose that \(G\) is an extension of an
abelian variety by a linear group \(L\) that is smooth over
\(k\).  Then \(L_K\) necessarily contains \(R\), so \(L=G\).
But \(G\) cannot be a linear group, because \(G_K\) has a nonzero
abelian variety as a quotient.

\smallskip
\noindent\textit{(c)} Let \(S\) be a solvable radical of \(G\),
and let \(B\) be a Borel subgroup of \(G\).  Then \(S_K\) and
\(B_K\) contain \(R\), so \(S=G\) and \(B=G\).  But then \(G\)
is solvable, contradicting the fact that \(G_K\) has a quotient
that is not solvable.

\SGAProof{Proof of \SGARefLink{sga3:xvii:proposition:c-5-1}{Proposition C.5.1}}
\textit{\SGAInlineRefTarget{sga3:xvii:corollary:C:5:1:item:i}(i)} We begin by describing the morphism \(u\).  A
\(K\)-scheme
\[
  f:S\longrightarrow\operatorname{Spec}(K)
\]
gives rise to the diagram
\[
\begin{tikzcd}[column sep=6.0em,row sep=4.0em]
  \operatorname{Spec}(K)
    \arrow[d,"j"'] &
  T
    \arrow[l,"h"']
    \arrow[d,bend right=28,"j_T"'] \\
  \operatorname{Spec}(k) &
  S
    \arrow[l,"g"]
    \arrow[ul,"f"]
    \arrow[u,bend right=28,"s"']
\end{tikzcd}
\]
where
\[
  g=j\circ f,\qquad
  T=\operatorname{Spec}(K)\times_{\operatorname{Spec}(k)}S,
\]
\(h\) and \(j_T\) are the two projections, and \(s\) is the
section of \(T\) over \(S\) such that \(h\circ s=f\).  The map
\[
  u(S):G_K(S)\longrightarrow H(S)
\]
is simply the composite
\[
  G_K(S)\xrightarrow{\sim}H(T)\longrightarrow H(S),
\]
where the last arrow is defined by the section \(s\).

Take, in particular, \(S\) to be the spectrum of an algebraic
closure \(\bar k\) of \(k\), and let \(f\) be the unique
\(k\)-morphism induced by \(K\to\bar k\), so that \(T\) is a
local Artinian scheme.  To prove \SGARefLink{sga3:xvii:corollary:C:5:1:item:i}{(i)}, it is enough to do so after
extending the base field from \(K\) to \(\bar k\).  It is clear
that
\[
  G_S=\prod_{T/S}H_T/T
\]
represents the Greenberg functor of \(H_T\) relative to \(S\)
(M. J. Greenberg, ``Schemata over local rings,'' \emph{Ann. of
Math.} 73 (1961), 624--648).  Under this identification, the
description above shows that \(u_S\) is the canonical transition
morphism
\[
  \operatorname{Green}(H_T)\longrightarrow
  H_S=H_T\times_T S.
\]
Assertion \SGARefLink{sga3:xvii:corollary:C:5:1:item:i}{(i)} now follows from the smoothness of \(H\) over \(K\)
and M. J. Greenberg, ``Schemata over local rings II,''
\emph{Ann. of Math.} 78 (1963), 256--266.

\smallskip
\noindent\textit{\SGAInlineRefTarget{sga3:xvii:corollary:C:5:1:item:ii}(ii)} To prove \SGARefLink{sga3:xvii:corollary:C:5:1:item:ii}{(ii)}, we may suppose that \(k\)
is separably closed.  Let \(U\) be a smooth algebraic subgroup of
\(G\) such that \(U_K\supset R\), and let us prove that \(U=G\).
Since \(G\) is connected, we may suppose that \(U\) is connected.
Let \(V=u(U_K)\), a smooth, connected algebraic subgroup of
\(H\), and put
\[
  V'=\prod_{K/k}V/K,
\]
a smooth, connected algebraic subgroup of \(G\).  The group
\(V'_K\) is an algebraic subgroup of \(G_K\), and the canonical
morphism \(V'_K\to V\) is the restriction of \(u\) to \(V'_K\).
By hypothesis, \(U_K\) contains \(R\), hence \emph{a fortiori}
contains \(\ker(V'_K\to V)\); by construction, moreover, the
image of \(U_K\) in \(H\) is \(V\).  It follows that \(U_K\)
contains \(V'_K\), and hence that \(U\) contains \(V'\).

On the other hand, the canonical isomorphism
\[
  G(k)\lhook\joinrel\longrightarrow G(K)
    \xrightarrow{u(K)}H(K)
\]
plainly sends \(U(k)\) into \(V(K)=V'(k)\).  Thus
\(U(k)\subset V'(k)\).  Since \(U\) is smooth and \(k\) is
separably closed, it follows that \(U\subset V'\), whence
\(U=V'\).  We then have
\[
  ([K:k]-1)\dim V
  =\dim\ker(V'_K\to V)
  =\dim R
  =([K:k]-1)\dim H.
\]
Since \(K\ne k\), we conclude that \(\dim V=\dim H\), hence that
\(V=H\), and finally that \(U=V'=G\).
